\section*{Теория}
\addcontentsline{toc}{section}{Теория}

\subsection*{1. Билинейные и квадратичные формы}

\begin{definition}
\textbf{Билинейной формой} на вещественном векторном пространстве $V$ называется функция $B: V \times V \to \mathbb{R}$, линейная по каждому из двух аргументов. 
В фиксированном базисе $\mathcal{E} = (e_1, \dots, e_n)$ билинейная форма задается своей матрицей $B$, элементы которой вычисляются как $b_{ij} = B(e_i, e_j)$. Значение формы на векторах $x$ и $y$ вычисляется по формуле:
\[
B(x,y) = [x]_\mathcal{E}^T \cdot B \cdot [y]_\mathcal{E}
\]
\end{definition}

\begin{definition}
\textbf{Квадратичной формой} $Q(x)$ на векторном пространстве $V$ называется функция $Q: V \to \mathbb{R}$, порожденная некоторой билинейной формой $B(x,y)$ при совпадении аргументов: $Q(x) = B(x,x)$.
\end{definition}

\textbf{Формула поляризации:} 
Любая квадратичная форма $Q(x)$ однозначно определяет порождающую её симметричную билинейную форму $B(x,y)$. Восстановить билинейную форму можно по формуле:
\[
B(x,y) = \frac{1}{2}\big( Q(x+y) - Q(x) - Q(y) \big)
\]

\textbf{Важное замечание (Выделение симметрической части):} 
Любая квадратная матрица $B$ может быть представлена в виде суммы симметрической и кососимметрической матриц: $B = \frac{B+B^T}{2} + \frac{B-B^T}{2}$.
Поскольку для любого вектора $x$ выполняется $x^T \frac{B-B^T}{2} x = 0$, значение квадратичной формы определяется \textbf{только симметрической частью} исходной матрицы.
Поэтому матрица квадратичной формы $\hat{A}$ по определению задается как:
\[
\hat{A} = \frac{B + B^T}{2}
\]
где $\hat{A}^T = \hat{A}$.

\subsection*{2. Смена базиса для билинейной и квадратичной форм}

Пусть задан переход от старого базиса $\mathcal{E}$ к нового базису $\mathcal{\hat{E}}$ с помощью невырожденной матрицы перехода $C$ (где столбцы $C$ — координаты векторов нового базиса в старом). 
Тогда координаты вектора преобразуются как $[x]_\mathcal{E} = C [x]_\mathcal{\hat{E}}$.

При замене базиса матрицы \textbf{билинейной формы} $B$ и \textbf{квадратичной формы} $A$ преобразуются конгруэнтно по следующим формулам:
\[
\hat{B} = C^T B C, \qquad \hat{A} = C^T A C
\]
\textit{Внимание: закон преобразования форм ($C^T A C$) кардинально отличается от закона преобразования линейных операторов ($C^{-1} A C$). След и определитель матрицы формы НЕ сохраняются при смене базиса (сохраняется лишь знак определителя).}

\subsection*{3. Приведение к каноническому и нормальному виду}

\begin{definition}
\textbf{Каноническим видом} квадратичной формы называется такой базис, в котором её матрица диагональна, то есть форма не содержит попарных произведений переменных:
\[
Q(y) = \lambda_1 y_1^2 + \lambda_2 y_2^2 + \dots + \lambda_n y_n^2
\]
Если коэффициенты $\lambda_i$ могут принимать только значения из множества $\{-1, 0, 1\}$, такой вид называется \textbf{нормальным}.
\end{definition}

\textbf{Чтение характеристик по каноническому виду:}
Пусть форма приведена к каноническому виду.
\begin{itemize}
    \item \textbf{Положительный индекс инерции ($p$):} количество строго положительных коэффициентов.
    \item \textbf{Отрицательный индекс инерции ($q$):} количество строго отрицательных коэффициентов.
    \item \textbf{Ранг формы ($r$):} число ненулевых коэффициентов ($r = p + q$).
    \item \textbf{Сигнатура:} пара чисел $(p, q)$ или разность $p - q$.
\end{itemize}
Согласно \textbf{закону инерции Сильвестра}, индексы инерции $p$ и $q$ зависят только от самой квадратичной формы и не зависят от способа ее приведения.

\subsection*{4. Приведение к каноническому виду: Метод Лагранжа и Симметричный Гаусс}

\textbf{Метод Лагранжа (алгебраический):} 
Алгоритмический процесс приведения формы к каноническому виду путем последовательного выделения полных квадратов.
\begin{enumerate}
    \item Если есть член с квадратом ($a_{11} x_1^2 \neq 0$), собираем все слагаемые с $x_1$ и дополняем их до полного квадрата.
    \item Если квадратов нет, но есть перекрестное произведение ($a_{12} x_1 x_2 \neq 0$), делаем невырожденную замену: $x_1 = y_1 + y_2$, $x_2 = y_1 - y_2$, чтобы искусственно создать квадраты, после чего возвращаемся к шагу 1.
\end{enumerate}

\textbf{Симметричный метод Гаусса (матричный):}
Матричный эквивалент метода Лагранжа. Позволяет одновременно найти диагональный вид матрицы и матрицу перехода $C$.
\begin{enumerate}
    \item Выписываем блочную матрицу $(A \mid E)$.
    \item Выполняем элементарное преобразование над \textbf{строками} матрицы $A$, и это же преобразование применяем к строкам матрицы $E$.
    \item Сразу же выполняем \textbf{точно такое же} элементарное преобразование над \textbf{столбцами} матрицы $A$ (при этом матрицу $E$ по столбцам НЕ трогаем!).
    \item Повторяем шаги 2-3, пока на месте $A$ не образуется диагональная матрица $D$.
    \item Итоговый вид блочной матрицы: $(D \mid C^T)$. Транспонировав правую часть, получаем искомую матрицу перехода $C$.
\end{enumerate}

\subsection*{5. Ортогональное приведение (Спектральная теорема)}

Любую вещественную симметрическую матрицу $A$ можно привести к диагональному виду с помощью \textit{ортогонального преобразования координат} (матрица перехода $C$ является ортогональной, т.е. $C^T = C^{-1}$). В этом случае:
\[
\hat{A} = C^T A C = C^{-1} A C = \operatorname{diag}(\lambda_1, \dots, \lambda_n)
\]
При таком преобразовании коэффициенты канонического вида в точности совпадают с \textbf{собственными значениями} матрицы $A$, а столбцы ортогональной матрицы $C$ образуют ортонормированный базис из \textbf{собственных векторов}.

\subsection*{6. Знакоопределенность и Критерий Сильвестра}

Квадратичная форма $Q(x)$ называется:
\begin{itemize}
    \item \textbf{Положительно определенной} ($Q(x) > 0$ при $x \neq 0$);
    \item \textbf{Отрицательно определенной} ($Q(x) < 0$ при $x \neq 0$);
    \item \textbf{Неотрицательно определенной / полуопределенной} ($Q(x) \ge 0$ для всех $x$);
    \item \textbf{Неположительно определенной / полуопределенной} ($Q(x) \le 0$ для всех $x$).
\end{itemize}

\begin{theorem}[Критерий Сильвестра]
Пусть $A$ — симметрическая матрица квадратичной формы, а $\Delta_1, \Delta_2, \dots, \Delta_n$ — её \textbf{главные угловые} миноры.
\begin{enumerate}
    \item Форма \textbf{строго положительно определена} $\iff$ $\Delta_1 > 0, \Delta_2 > 0, \dots, \Delta_n > 0$.
    \item Форма \textbf{строго отрицательно определена} $\iff$ знаки чередуются, начиная с минуса: $\Delta_1 < 0, \Delta_2 > 0, \Delta_3 < 0, \dots$
\end{enumerate}
\end{theorem}

\textbf{Обобщенный критерий:} 
Для того чтобы форма была \textbf{неотрицательно (нестрого) определенной}, \underline{недостаточно} условия $\Delta_k \ge 0$ только для угловых миноров. Необходимо и достаточно, чтобы \textbf{абсолютно все главные миноры} матрицы (всех порядков, а не только стоящие в левом верхнем углу) были $\ge 0$.
